\begin{resumo}[Abstract]

    \begin{otherlanguage*}{english}

        \noindent 
            The expansion of the internet and communication media has increased access to content in different languages, such as videos, articles, and blogs. In this context, machine translation tools contribute to reducing language barriers by enabling the conversion of information between different languages. This study compares three neural machine translation models based on the Transformer architecture and available with open weights: MarianMT, M2M-100, and NLLB. The models were evaluated in English-to-Portuguese translation, considering both general and domain-specific texts. Translation quality was analyzed using the BLEU, BERTScore, and chrF metrics, while computational efficiency was assessed based on average CPU usage, inference time, and peak RAM consumption. Under the established experimental conditions, MarianMT achieved the best results across all three quality metrics, as well as the lowest average CPU usage and inference time. Although differences in processing performance were observed among the models, their peak RAM consumption values were similar. It was also found that domain-specific texts resulted in lower scores than general texts across the quality metrics. Therefore, within the analyzed scenario, MarianMT demonstrated the best trade-off between translation quality and computational efficiency, while peak RAM consumption did not allow for significant differentiation among the models.
        \noindent  
        
        \textbf{Keywords}: transformer, machine translation, automatic metrics, computational efficiency.
 
    \end{otherlanguage*}

\end{resumo}